\itemhead{习题4.1 · 第10题（教材p.105）}{XJ-B1-C04-HW-P105-Y0405}
\label{XJ-B1-C04-HW-P105-Y0405}
{\footnotesize 湘教版教材p.105 / PDF第112页 / 习题4.1 / 第10题}
\field{原题}{已知 $x>1$，且 $x+x^{-1}=3$，求下列各式的值：
\subq{1}$x^{1/2}+x^{-1/2}$；
\subq{2}$x^{1/2}-x^{-1/2}$；
\subq{3}$x^{3/2}+x^{-3/2}$。}
\field{参考答案}{依次为 $\sqrt5,1,2\sqrt5$。第二问由 $x>1$ 选正号；第三问用 $(u+v)^3-3uv(u+v)$，其中 $uv=1$。}
{\small \field{关联知识}{P02 有理数、实数指数幂与运算（教材帮p.144、145 / PDF153、154）}
\field{方法与易错点}{统一为指数幂，分开处理系数与各底数的指数；最终复查定义条件及数值精度。 易错：负指数误作负数；幂的幂指数相加；约分丢禁点；过早四舍五入。}
\field{答案依据}{答案册相关题组：PDF第31页（册内页30）；按当前题库原题逐项重算，不按旧版题号机械对套。}}
